\section{Methodology}
\label{sec:methodology}

\subsection{Problem Formulation}

We formulate Stage 0 constellation seeding as a sequential decision-making problem. At each step $t$, the agent observes the current constellation state and selects orbital elements for the next satellite. The episode terminates when the constellation achieves feasibility or reaches the maximum number of satellites.

\textbf{State Space:} The state consists of:
\begin{itemize}
    \item Placed satellites: $\mathcal{S}_t = \{s_1, \dots, s_{n_t}\}$ where each $s_i \in \mathbb{R}^6$ represents orbital elements $[a, e, i, \Omega, \omega, \nu]$.
    \item Debris clients: $\mathcal{C} = \{c_1, \dots, c_m\}$ where each $c_j \in \mathbb{R}^6$ represents client orbital elements.
    \item Orbital bounds: $B \in \mathbb{R}^{6 \times 2}$ defining feasible ranges for each orbital element.
    \item Remaining budget: $b_t = N_{\max} - n_t$ satellites remaining to place.
    \item Cost state: Current deficit, unsafe cost, safe cost, and feasibility status.
\end{itemize}

\textbf{Action Space:} The action is a 6-dimensional vector representing the orbital elements of the next satellite: $a_t \in \mathbb{R}^6$. Actions are normalized to $[-1, 1]$ and denormalized to orbital bounds during execution.

\textbf{Reward Function:} The reward at step $t$ is:
\begin{equation}
    r_t = \Delta C_t - \lambda_u C_u - \lambda_s C_s - \lambda_p C_p + \mathbb{I}_{\text{feasible}} \cdot R_{\text{bonus}}
\end{equation}
where $\Delta C_t$ is the cost improvement, $C_u, C_s, C_p$ are unsafe, safe, and prediction costs, $\lambda_u, \lambda_s, \lambda_p$ are penalty weights, and $R_{\text{bonus}}$ is a bonus for achieving feasibility.

\subsection{Environment Design}

The \texttt{SatelliteSeedingEnv} implements the RL environment following the OpenAI Gym interface. Key features include:

\begin{itemize}
    \item \textbf{Reset:} Initialize with empty constellation and load client data.
    \item \textbf{Step:} Place satellite, compute cost, update state, check termination.
    \item \textbf{Cost Function:} Supports both simplified orbital distance proxy and full CAPO audit.
\end{itemize}

The environment uses the existing \texttt{compute\_stage0\_cost} function from the SpaceAGORA framework, ensuring consistency with the baseline methods.

\subsection{Policy Architecture}

We use a dual DeepSet architecture (Figure~\ref{fig:architecture}) with the following components:

\begin{enumerate}
    \item \textbf{Satellite Encoder:} $\phi_s: \mathbb{R}^6 \rightarrow \mathbb{R}^{d_{\text{embed}}}$ encodes each satellite.
    \item \textbf{Client Encoder:} $\phi_c: \mathbb{R}^6 \rightarrow \mathbb{R}^{d_{\text{embed}}}$ encodes each debris client.
    \item \textbf{Aggregation:} Sum embeddings for each set: $h_s = \sum_{s \in \mathcal{S}} \phi_s(s)$, $h_c = \sum_{c \in \mathcal{C}} \phi_c(c)$.
    \item \textbf{Cross-Attention:} $h = \text{Attention}(h_s, h_c)$ captures interactions.
    \item \textbf{Policy Head:} $\pi: \mathbb{R}^{d_{\text{hidden}}} \rightarrow \mathbb{R}^6$ outputs normalized orbital elements.
\end{enumerate}

The policy is trained using PPO with the following hyperparameters:
\begin{itemize}
    \item Learning rate: $3 \times 10^{-4}$
    \item Hidden size: 32
    \item Embedding size: 16
    \item PPO clip: 0.2
    \item GAE $\lambda$: 0.95
    \item Discount $\gamma$: 0.99
\end{itemize}

\subsection{Training Infrastructure}

We implement a comprehensive training campaign infrastructure that supports:

\begin{itemize}
    \item \textbf{Cluster Combinations:} Individual clusters (1a, 2a, 3a...), pairs (1a+2a, 2a+3a...), and supersets (1a+2a+3a...).
    \item \textbf{Parameter Sweeps:} Full-factorial sampling over laser range, laser power, and planning horizons.
    \item \textbf{Curriculum Learning:} Sort scenarios by difficulty (easy → hard) for progressive training.
    \item \textbf{Checkpoint/Resume:} Compatible with existing PPOTrainer checkpoint system using JLD2 format.
    \item \textbf{TensorBoard Integration:} Log training metrics, embeddings, and cost components for visualization.
\end{itemize}

The training entry point \texttt{run\_rl\_training\_campaign.jl} parses a YAML configuration, generates scenarios, and executes training in parallel or sequential mode.

\subsection{ROE-Based Orbital Bounds}

To maintain compatibility with CAPOConstellation without external dependencies, we port the ROE math natively to SpaceAGORA. The \texttt{ROEBoundsCalculator} module provides:

\begin{itemize}
    \item \texttt{compute\_orbital\_bounds\_from\_cluster:} Statistical bounds from debris cluster orbital elements.
    \item \texttt{roe\_to\_orbital\_elements:} Convert ROE to absolute orbital elements.
    \item \texttt{orbital\_elements\_to\_roe:} Convert absolute orbital elements to ROE.
    \item \texttt{compute\_reference\_orbit:} Compute reference orbit from cluster mean.
\end{itemize}

This ensures that the RL search space bounds match those used by stochastic greedy seeding for fair comparison.
